\section{Integrantes y áreas de trabajo del \texorpdfstring{\ac{GRID}}{GRID}}

El Grupo de Investigación en Redes, Información y Distribución (\ac{GRID}) está conformado por docentes e investigadores vinculados a diferentes áreas de conocimiento relacionadas con las tecnologías de la información y las comunicaciones. Sus integrantes participan en actividades de investigación, desarrollo y formación, haciendo uso de los recursos y servicios tecnológicos disponibles en la infraestructura del grupo.

A continuación, se presenta una lista de los integrantes actuales del grupo, junto
con sus respectivas líneas de trabajo y áreas de especialización:

\begin{itemize}
	\item \href{https://scienti.minciencias.gov.co/cvlac/visualizador/generarCurriculoCv.do?cod_rh=0000210897}{\textbf{Christian Andrés Candela Uribe}}: Microservicios, desarrollo de software, minería de datos, infraestructura \ac{TI}.
	\item \href{https://scienti.minciencias.gov.co/cvlac/visualizador/generarCurriculoCv.do?cod_rh=0001383939}{\textbf{Luis Eduardo Sepúlveda Rodríguez}}: Infraestructura \ac{TI}, \ac{HPC}, computación paralela.
	\item \href{https://scienti.minciencias.gov.co/cvlac/visualizador/generarCurriculoCv.do?cod_rh=0001343801}{\textbf{Carlos Eduardo Gómez Montoya}}: Redes, ingeniería de software, cloud computing.
	\item \href{https://scienti.minciencias.gov.co/cvlac/visualizador/generarCurriculoCv.do?cod_rh=0001398775}{\textbf{Sergio Augusto Cardona Torres}}: Big data y análisis de datos, ingeniería de software, sistemas adaptativos, informática educativa.
	\item \href{https://scienti.minciencias.gov.co/cvlac/visualizador/generarCurriculoCv.do?cod_rh=0000193550}{\textbf{Sonia Jaramillo Valbuena}}: Big data, electroquímica, inteligencia artificial.
	\item \href{https://scienti.minciencias.gov.co/cvlac/visualizador/generarCurriculoCv.do?cod_rh=0000283495}{\textbf{Julián Esteban Gutiérrez Posada}}: Microservicios, desarrollo de software, minería de datos, infraestructura \ac{TI}, \ac{HPC}, computación paralela, redes, ingeniería de software.
\end{itemize}

La diversidad de áreas de trabajo de los integrantes del \ac{GRID} aporta capacidades complementarias en infraestructura de \ac{TI}, redes, desarrollo de software y gestión de servicios tecnológicos. Esta combinación de conocimientos favorece la atención de las necesidades técnicas asociadas con la infraestructura del grupo y proporciona las capacidades necesarias para integrar y administrar los diferentes servicios que la conforman. En este contexto, la experiencia de sus integrantes constituye un elemento relevante para la implementación de mecanismos orientados a centralizar la gestión de identidades y facilitar el control de acceso a los recursos tecnológicos disponibles.
